\chapter[Band 45: Halbverbände und Verbände]{Band 45 auf einen Blick}
\noindent\textbf{Halbverbände und Verbände}\par\medskip
Dieser Band stellt Halbverbände eigenständig axiomatisch dar und
verbindet sie mit partiellen Ordnungen. Wir schreiben die induzierte
Supremumsordnung kurz als \(\le\). Endlichkeit und ein Nullelement
werden erst bei den entsprechend bezeichneten Resultaten vorausgesetzt.

\begin{BandUebersicht}
\textbf{Halbverbandsaxiome}
Eine Halbverbandsoperation \(\star:A\times A\to A\) erfüllt für
\(x,y,z\in A\):\UebFormel{\begin{aligned}
&(x\star y)\star z=x\star(y\star z),\\
&x\star y=y\star x,\qquad x\star x=x.\\
&x\le y\ \leftrightarrow\
 (x,y\in A\land x\star y=y).
\end{aligned}}
\UebQuellen{\FormulaRefAuto{Semilattice}[def];
\FormulaRefAuto{JoinOrder}[def]}
&
\textbf{Partielle Ordnung und Paarsupremum}
\UebFormel{\begin{aligned}
&\Semi{A,\star}\ \vdash\PartOrd{A,\le},\\
&\Semi{A,\star},\ x,y\in A\ \vdash\\
&\qquad x\star y=\sup_{A,\le}\{x,y\}.
\end{aligned}}
\UebQuellen{\FormulaRefAuto{JoinSemilatticeInducesOrder}[thm];
\FormulaRefAuto{JoinIsPairSupremum}[thm]} \\
\addlinespace[.8ex]

\textbf{Morphismen von Halbverbänden}
Für Halbverbände \((A,\star)\), \((B,\diamond)\) ist
\(f:A\to B\) ein Homomorphismus, wenn\UebFormel{\begin{aligned}
&x,y\in A\ \vdash\\
&\quad f(x\star y)=f(x)\diamond f(y).
\end{aligned}}Ein Isomorphismus ist ein bijektiver Homomorphismus.
\UebQuellen{\FormulaRefAuto{SemilatticeHomomorphism}[def];
\FormulaRefAuto{SemilatticeIsomorphism}[def]}
&
\textbf{Komposition}
Sind \(f:A\to B\) und \(g:B\to C\) Halbverbandshomomorphismen,
dann ist auch \(g\circ f:A\to C\) einer; insbesondere gilt\UebFormel{\begin{aligned}
&(g\circ f)(x\star y)\\
&\quad=(g\circ f)(x)\mathbin{\bullet}(g\circ f)(y).
\end{aligned}}Für Isomorphismen gilt dieselbe Abgeschlossenheit.
\UebQuellen{\FormulaRefAuto{SemilatticeHomComposition}[thm];
\FormulaRefAuto{SemilatticeIsoComposition}[thm]} \\
\addlinespace[.8ex]

\textbf{Vom Paarsupremum zur Algebra}
In einer partiell geordneten Menge \((A,\le)\) mit allen
Paarsuprema wird definiert:\UebFormel{\begin{aligned}
&x\star y=\sup_{A,\le}\{x,y\}
 \quad(x,y\in A).
\end{aligned}}
\UebQuellen{\FormulaRefAuto{PairJoinOperation}[def]}
&
\textbf{Rückgewinnung der Ordnung}
Diese Operation ist eine Halbverbandsoperation und gewinnt
die Ausgangsordnung zurück:\UebFormel{\begin{aligned}
&\Semi{A,\star},\\
&x,y\in A\ \vdash\
 (x\le y\leftrightarrow x\star y=y).
\end{aligned}}
\UebQuellen{\FormulaRefAuto{PairJoinCharacterization}[thm];
\FormulaRefAuto{PairJoinRecoversOrder}[thm]} \\
\addlinespace[.8ex]

\textbf{Infimale Lesart und Mengenbeispiele}
Infimal liest man eine Halbverbandsoperation durch\UebFormel{\begin{aligned}
&x\le_{\mathrm{inf}}y
 \ \leftrightarrow x\star y=x
 \quad(x,y\in A).
\end{aligned}}Auf \(\powerset(M)\) stehen dafür \(\cap\) und \(\cup\) zur Verfügung.
\UebQuellen{\FormulaRefAuto{MeetOrderEvaluation}[thm];
\FormulaRefAuto{PowerSetUnionSemilattice}[thm];
\FormulaRefAuto{PowerSetIntersectionSemilattice}[thm]}
&
\textbf{Dualität und Inklusion}
Bei der infimalen Lesart liefert \(x\star y\) das Paarinfimum.
Auf der Potenzmenge gilt\UebFormel{\begin{aligned}
&X\cup Y=Y\ \leftrightarrow X\subseteq Y,\\
&X\cap Y=X\ \leftrightarrow X\subseteq Y\\
&\hfill(X,Y\subseteq M).
\end{aligned}}
\UebQuellen{\FormulaRefAuto{MeetIsPairInfimum}[thm];
\FormulaRefAuto{PowerSetUnionOrderIsInclusion}[thm];
\FormulaRefAuto{PowerSetIntersectionOrderIsInclusion}[thm]} \\
\addlinespace[.8ex]

\textbf{Verband und Verbandshomomorphismus}
Ein Verband \((A,\MeetOp,\JoinOp)\) besitzt zwei
Halbverbandsoperationen und erfüllt\UebFormel{\begin{aligned}
&x\MeetOp(x\JoinOp y)=x,\\
&x\JoinOp(x\MeetOp y)=x\quad(x,y\in A).
\end{aligned}}Ein Verbandshomomorphismus erhält beide Operationen.
\UebQuellen{\FormulaRefAuto{AlgebraicLattice}[def];
\FormulaRefAuto{LatticeHomomorphism}[def];
\FormulaRefAuto{LatticeIsomorphism}[def]}
&
\textbf{Eine gemeinsame Ordnung}
\UebFormel{\begin{aligned}
&\LatticeAlg{A,\MeetOp,\JoinOp},\ x,y\in A\\
&\vdash (x\MeetOp y=x\leftrightarrow x\JoinOp y=y),\\
&\LatticeAlg{\powerset(M),\cap,\cup}.
\end{aligned}}Verbandshomomorphismen sind unter Komposition abgeschlossen.
\UebQuellen{\FormulaRefAuto{LatticeOrdersCoincide}[thm];
\FormulaRefAuto{PowerSetLattice}[thm];
\FormulaRefAuto{LatticeHomComposition}[thm]} \\
\addlinespace[.8ex]

\textbf{Null und Eins}
Ein Supremums-Halbverband mit Null besitzt \(0\in A\) mit
\(0\JoinOp x=x\). Ein beschränkter besitzt zusätzlich \(1\in A\) mit
\(x\JoinOp1=1\), jeweils für \(x\in A\).
\UebQuellen{\FormulaRefAuto{ZeroJoinSemilattice}[def];
\FormulaRefAuto{BoundedJoinSemilattice}[def]}
&
\textbf{Schranken und endliche Vollständigkeit}
Null ist das kleinste und Eins das größte Element.
Jeder endliche Supremums-Halbverband mit Null besitzt ein größtes Element:\UebFormel{\begin{aligned}
&\Finite A,\ \ZeroJoinSemi{A,\JoinOp,0}\ \vdash\\
&\quad\exists t\in A\ \forall x\in A\ (x\le t).
\end{aligned}}
\UebQuellen{\FormulaRefAuto{ZeroJoinLeast}[thm];
\FormulaRefAuto{BoundedJoinTopGreatest}[thm];
\FormulaRefAuto{FiniteZeroJoinSemilatticeHasGreatest}[thm];
\FormulaRefAuto{FiniteZeroJoinSemilatticeIsBounded}[thm]} \\
\addlinespace[.8ex]

\textbf{Gemeinsame untere Schranken}
Für \(x,y\in A\) setze\UebFormel{\begin{aligned}
&D(x,y)=\{d\in A\mid d\le x,\ d\le y\}.
\end{aligned}}In einem Supremums-Halbverband mit Null ist \(0\in D(x,y)\).
\UebQuellen{\FormulaRefAuto{ZeroJoinLeast}[thm]}
&
\textbf{Endliche Halbverbände mit Null haben Paarinfima}
\UebFormel{\begin{aligned}
&\Finite A,\ \ZeroJoinSemi{A,\JoinOp,0},\\
&x,y\in A\ \vdash\\
&\exists m\in D(x,y)\ \forall d\in D(x,y)\
 (d\le m).
\end{aligned}}Dieses größte untere Element ist das Paarinfimum.
\UebQuellen{\FormulaRefAuto{FiniteZeroJoinSemilatticePairInfima}[thm]} \\
\addlinespace[.8ex]

\textbf{Irreduzible und Trenner}
In einem Halbverband mit Null heißt \(j\in A\setminus\{0\}\)
\(\JoinOp\)-irreduzibel, wenn\UebFormel{\begin{aligned}
&x,y\in A,\quad x\JoinOp y=j\\
&\qquad\vdash x=j\lor y=j.\\
&\operatorname{Sep}(x,y)
 =\{d\in A\mid d\le x,\ d\not\le y\}.
\end{aligned}}
\UebQuellen{\FormulaRefAuto{JoinIrreducible}[def];
\FormulaRefAuto{JoinSeparatorSet}[def]}
&
\textbf{Irreduzible trennen}
Minimale Trenner sind irreduzibel. Ist außerdem \(A\) endlich,
so gilt für \(x,y\in A\):\UebFormel{\begin{aligned}
&x\not\le y\ \vdash\\
&\exists j\in\JoinIrrSet{A,\JoinOp,0}\\
&\qquad(j\le x\land j\not\le y).
\end{aligned}}
\UebQuellen{\FormulaRefAuto{MinimalSeparatorJoinIrreducible}[thm];
\FormulaRefAuto{JoinIrreduciblesSeparate}[thm]} \\
\addlinespace[.8ex]

\textbf{Echte untere Elemente}
In der im Band verwendeten Notation für einen Halbverband \(L\)
mit Nullelement \(\varnothing\):\UebFormel{\begin{aligned}
&D_L(P)=\\
&\quad\{Q\in L\mid Q\le P,\ Q\ne P\}.
\end{aligned}}
\UebQuellen{\FormulaRefAuto{JoinIrreducibleProperLowerSet}[def]}
&
\textbf{Größtes echtes unteres Element}
\UebFormel{\begin{aligned}
&\Finite L,\ \ZeroJoinSemi{L,\JoinOp,\varnothing},\\
&\JoinIrr{L,\JoinOp,\varnothing}P\ \vdash\\
&\exists P_\ast\in D_L(P)\
 \forall Q\in D_L(P)\ (Q\le P_\ast).
\end{aligned}}Diese Eigenschaft trägt später den Beweis über private Punkte.
\UebQuellen{\FormulaRefAuto{JoinIrreducibleLargestProperLower}[thm]} \\
\addlinespace[.8ex]
\end{BandUebersicht}

